\newpage
\section{Conclusão}
\label{sec:conclusao}

O objetivo deste trabalho foi projetar e validar um rastreador de referência para um modelo farmacocinético discreto, considerando incertezas paramétricas na dinâmica da planta. Para isso, foi escolhida uma malha de controle PI, e através da sua representação em espaço de estados, as condições necessárias para sua estabilidade foram determinadas e expressas por meio de LMIs. Os ganhos do controlador foram sintetizados através de um procedimento iterativo baseado em tais condições, onde as incertezas paramétricas esperadas foram modeladas mediante um polítopo de oito vertices.

Uma primeira avaliação computacional foi realizada em Matlab e Simulink, utilizando os mesmos ganhos determinados pelo procedimento iterativo. Nas simulações realizadas, a saída da planta se aproximou da referência em todos os vértices do polítopo, sem sobressinal. No entanto, dependendo do valor dos parâmetros incertos, observaram-se diferenças no tempo de estabilização do sinal de saída.

Para validar tais resultados, por meio de dois microcontroladores ESP32 foi realizada uma emulação embarcada do sistema. A comunicação entre os dois dispositivos, um ESP32 planta e um ESP32 controlador, foi estabelecida por UART, com uma interface gráfica desenvolvida em Node-RED mediante o protocolo MQTT  para registrar o comportamento do sistema e alterar a referência. Para representar os parâmetros incertos, três potenciômetros foram conectados ao ESP32 planta, permitindo percorrer trajetorias dentro do volume descrito pelo polítopo. Nos ensaios com esta plataforma, foi avaliado o comportamento do sistema nos vértices V4 e V5 para dois degraus, um de subida e outro de descida. Em ambos os casos a saída conseguiu acompanhar a referência desejada, mas as respostas transitórias apresentaram diferenças. A saida no vértice V4 não apresentou sobressinal, enquanto no vértice V5 a saída ultrapassou temporariamente a referência, com um sobressinal estimado em $\approx6,67\%$. Esse comportamento foi exclusivo à emulação embarcada. Além dos testes com parâmetros fixos, foram realizados ensaios variando os três parâmetros incertos durante o funcionamento da malha fechada, observado que mesmo diante variação, a malha permaneceu estável.

Dessa forma, os objetivos de projeto do controlador, simulação da malha fechada e emulação embarcada foram atingidos. Os resultados experimentais mostram que a estabilidade foi preservada nas condições avaliadas, inclusive ao variar os parâmetros incertos dentro do politopo. No entanto, esses ensaios não representam uma verificação de estabilidade para todas as trajetórias possíveis de variação dos parâmetros. Devido ás diferenças entre a emulação embarcada e a simulação, propõe-se ampliar os testes realizados na plataforma embarcada, investigando a influência da temporização das tasks e da comunicação UART, com o objetivo de realizar uma comparação quantitativa entre a implementação computacional e o sistema embarcado.

